\chapter{Proofs for the fixed-priority Prakken--Sartor system}
\label{annex:p18-is1996}

\begin{annexitem}{annex:p18.is1996-closure}{thm:p18-is1996-closure}{Strict saturation and valid extensions}
\begin{proof}
Every conclusion of a strict extension is in the Horn closure
of the original conclusions: induction along the appended list
uses the defining strict rule at each step.
Conversely, take a finite strict derivation from
\(\operatorname{Con}(A)\).
Process it in derivation order. If a rule's consequent has already
been produced, omit that rule; otherwise append it.
Its strong premises have already been produced, either in \(A\)
or at an earlier processed step. Omitting a redundant producer
does not remove the literal that later steps need.
The resulting list has unique producers and is a strict extension
containing the desired conclusion. This proves equality.

Subargument inclusion gives inclusion of conclusions and of recorded
weak premises. Horn closure is monotone in its initial facts,
proving both asserted inclusions.
For a repeated identical rule, retain its first occurrence.
All later uses of its consequent still have that earlier producer,
and the identical weak premises were already recorded at the
first occurrence. Delete further occurrences successively.
This preserves conclusions, assumptions, and validity.
\end{proof}
\end{annexitem}

\begin{annexitem}{annex:p18.is1996-frontier}{thm:p18-is1996-frontier}{Uniqueness of relevant frontiers}
\begin{proof}
Each produced literal has exactly one producing rule.
If that rule is defeasible, the singleton frontier is specified.
If it is strict, every strong antecedent has a producer earlier
in the list. Recursion therefore descends in a finite ordering.
The empty-premise strict case gives the empty union.
Induction along the list proves existence, uniqueness and finiteness.
The same induction gives an empty frontier at every strict
conclusion of a wholly strict argument.
At a defeasible producer the recurrence stops, so earlier
defeasible rules do not enter that frontier. They remain in
the argument and their consequents and assumptions remain
available to the attack definition.
\end{proof}
\end{annexitem}

\begin{annexitem}{annex:p18.is1996-set-priority}{thm:p18-is1996-set-priority}{The strict lift, including the empty set}
\begin{proof}
If \(R=\varnothing\), the existential clause is false.
If \(R\ne\varnothing\) and \(R'=\varnothing\), any \(r\in R\)
is a witness because the universal clause is vacuous.
For irreflexivity, a proposed witness \(r\in R\) to \(R\prec_*R\)
would require \(r\prec r\), impossible.
Suppose \(R\prec_*R'\prec_*R''\), with witnesses
\(r\in R\) and \(r'\in R'\).
Then \(r\prec r'\), and \(r'\prec t\) for every \(t\in R''\).
Transitivity gives \(r\prec t\) for every such \(t\), so
\(R\prec_*R''\). This also covers \(R''=\varnothing\).
The intermediate set cannot be empty because it supplies a witness.
Transitivity and irreflexivity imply asymmetry.
\end{proof}
\end{annexitem}

\begin{annexitem}{annex:p18.is1996-orientation}{thm:p18-is1996-orientation}{Attack, defeat, and the special empty argument}
\begin{proof}
If either argument undercuts the other, that direction is a defeat
by definition. Otherwise take a head-conflict witness with frontiers
\(R,R'\). They cannot satisfy both \(R\prec_*R'\) and
\(R'\prec_*R\). Consequently at least one of the two rebuttal
conditions holds, and with no reverse undercut it is a defeat.
This proves orientation of every attack.
Rebuttal and undercutting retain their respective attack witnesses,
so each implies attack. The defeat disjunction adds only the
explicit empty/incoherent case, establishing the qualified converse.

An attack on a subargument remains an attack on its superargument:
use the closure and assumption inclusions already proved.
For the last assertion, conflicting strict arguments give
complementary literals in the strict closure of the empty argument.
That closure is contained in every \(C_S(A)\), because the
empty argument is a subargument of every argument.
Every \(A\) therefore attacks itself.

The exception in the converse is necessary. With \(S=\varnothing\)
and rules \(\Rightarrow a,\Rightarrow\neg a\), their combined
argument is incoherent. The empty argument defeats it through
the designated clause but has no strict-closure conclusion
and hence does not attack it.
\end{proof}
\end{annexitem}

\begin{annexitem}{annex:p18.is1996-semantics}{thm:p18-is1996-semantics}{The source fixed point and its structural properties}
\begin{proof}
The operator \(F_\Gamma\) is monotone because enlarging the
set of available strict defenders cannot remove a witness.
Its transfinite iteration from the empty set is increasing,
with unions at limit stages. It stabilizes on the set
\(\mathcal A_\Gamma\): otherwise choosing one new argument at
each strict increase would exceed that set's cardinality.
Induction shows that every iterate is contained in every fixed
point, so the stabilized set is the unique least fixed point.

Let \(U\) be the union of the finite iterates.
They are contained in \(J_\Gamma\).
If \(\Gamma\) is finitary and \(A\in F_\Gamma(U)\), there
are finitely many defeaters of \(A\): each is an attacker or
the one designated empty defeater. Choose a defending witness
for each. These finitely many witnesses occur together in
some finite iterate, so \(A\) occurs in its successor.
Thus \(F_\Gamma(U)\subseteq U\).
The reverse inclusion follows from the increasing iteration.
Hence \(U\) is a fixed point and equals the least one.

We establish a useful comparison without replacing the rule-generated
relation. Let \(H(X)\) require a defender that defeats, rather
than strictly defeats, every \(D_\Gamma\)-defeater.
Write \(J_H\) for its least fixed point.
Since \(F_\Gamma(X)\subseteq H(X)\), least-fixed-point
induction gives \(J_\Gamma\subseteq J_H\).
For the reverse inclusion, use first-entry ranks in the
transfinite iteration of \(H\).
If \(A\) enters at stage \(\alpha+1\), each defeater \(B\)
has a defender \(C\) of smaller rank.
Choose such a defender of least rank.
If \(B\) defeated \(C\), the earlier defense required to
admit \(C\) would give a defender of \(B\) of still smaller
rank. This is impossible. Thus \(C\) strictly defeats \(B\).
Transfinite induction now places every member of \(J_H\)
in \(J_\Gamma\). Therefore \(J_H=J_\Gamma\).

For completeness, the ordinary-defense iteration is conflict-free
under \(D_\Gamma\).
Call a set admissible if it is \(D_\Gamma\)-conflict-free
and contained in its image under \(H\).
If \(X\) is admissible and \(a,b\in H(X)\) with \(aD_\Gamma b\),
some \(c\in X\) defeats \(a\), and then some \(d\in X\)
defeats \(c\), contradicting conflict-freeness of \(X\).
Thus \(H(X)\) is conflict-free.
It defends itself because its defenses can be taken in
\(X\subseteq H(X)\).
Start with the empty admissible set. Unions of increasing
admissible chains are admissible: every pair lies together
at a stage, and each member retains its earlier defenders.
This proves conflict-freeness of \(J_H=J_\Gamma\).
Every attack between its members would yield a defeat in
one direction by the orientation theorem. Hence it is also
conflict-free under the source's attack relation.

Let \(A\in J_\Gamma\), let \(A'\) be a subargument, and
let \(B D_\Gamma A'\).
If this is the empty/incoherent case, incoherence of \(A'\)
implies incoherence of \(A\), so the same \(B\) defeats \(A\).
The fixed-point property supplies a strict defender in \(J_\Gamma\).
In every other case \(B\) attacks \(A'\), hence attacks \(A\).
If \(B\) defeats \(A\), again use its strict defender in
\(J_\Gamma\).
If it does not, orientation gives \(A D_\Gamma B\);
together with not \(B D_\Gamma A\), this is strict defeat.
Thus in every case the defeater \(B\) of \(A'\) is strictly
defeated by a member of \(J_\Gamma\).
Therefore \(A'\in F_\Gamma(J_\Gamma)=J_\Gamma\).

Finally, if the strict arguments are not conflict-free,
every argument is incoherent. The empty argument then
defeats every argument, including itself.
No argument can strictly defeat the empty argument,
since the empty argument defeats it in return.
Thus \(F_\Gamma(X)=\varnothing\) for every \(X\), and
\(J_\Gamma=\varnothing\).
The set-sized choices in this proof use ordinary classical
set theory with choice; the finite cases use finite choices only.
\end{proof}
\end{annexitem}

\begin{annexitem}{annex:p18.is1996-proof-correctness}{thm:p18-is1996-proof-correctness}{Finite proofs: soundness and finitary completeness}
\begin{proof}
For soundness, induct upward through a finite winning tree.
A proponent leaf has no defeaters and lies in
\(F_\Gamma(\varnothing)\).
At any other proponent node, every opponent defeater is
included among its children and has a proponent response
that strictly defeats it. The induction hypothesis puts
those responses in \(J_\Gamma\).
The node is therefore in \(F_\Gamma(J_\Gamma)=J_\Gamma\).
This gives the root claim, without any finitary assumption.

Now suppose \(\Gamma\) is finitary and \(A\in J_\Gamma\).
The approximation theorem places it in a finite iterate.
Assign each admitted argument its first finite entry stage.
For every opponent defeater of a proponent argument of rank \(n\),
choose a strict defender of rank less than \(n\), as provided
by the definition of that iterate.
Repeat the construction at each new proponent node.
Ranks strictly decrease at proponent turns, so no proponent
argument repeats and branch depth is bounded by twice the
root rank. Each proponent has finitely many defeaters,
and one response is chosen at each opponent node.
Induction on the depth bound shows the whole tree is finite.
Its leaves are proponent moves with no opponent reply,
giving the required proof.
No global minimality condition was imposed on the chosen defender.
\end{proof}
\end{annexitem}

\begin{annexitem}{annex:p18.is1996-minimality}{thm:p18-is1996-minimality}{The complete rule-generated minimality counterexample}
\begin{proof}
Take the four rules in the body, with \(S=\prec=\varnothing\).
Every subset of these rules is an argument content: there
are no strong premises, and their heads \(a,b,\neg a,\neg b\)
are distinct. Thus the universe has exactly sixteen contents,
including the empty one, and is finitary.

Let \(C=[r_1,r_2]\).
It has no assumptions and is coherent.
Any argument attacking \(C\) must conclude \(\neg a\) or
\(\neg b\), so must contain \(r_3\) or \(r_4\).
In the first case \(C\) undercuts its weak premise
\(\mathord{\sim}b\); in the second it undercuts
\(\mathord{\sim}a\).
The attacking argument cannot undercut \(C\), which has
no weak premises. Its rebuttal is therefore blocked by
the reverse undercut. The empty incoherence case does
not apply to coherent \(C\). Hence \(C\) has no defeaters.
It belongs to the first semantic iterate.
All defeaters of \([r_1]\) or \([r_2]\) are strictly
defeated by \(C\), by the same reasoning, so these
subarguments belong to the next iterate.

Consider a strict defeater \(Z\) of \([r_3]\).
If \(r_2\notin Z\), it cannot conclude \(b\), so it cannot
undercut \([r_3]\).
To defeat by rebuttal it must contain \(r_1\), producing \(a\).
But with no priorities, \([r_3]\) then rebuts \(Z\) in return;
that reverse rebuttal is not blocked by an undercut from \(Z\),
since \(Z\) cannot conclude \(b\).
So defeat cannot be strict.
Consequently every strict defeater of \([r_3]\) contains \(r_2\).
The singleton \([r_2]\) itself strictly defeats \([r_3]\):
it undercuts it and has no reverse defeater from that argument.
It is therefore the unique globally inclusion-minimal strict defender.
Symmetrically, \([r_1]\) is the unique globally minimal
strict defender of \([r_4]\).

Starting from \([r_1]\), opponent legally plays \([r_3]\).
Global minimality forces \([r_2]\), to which opponent
legally replies \([r_4]\).
The unique permitted minimal defense would repeat \([r_1]\).
Thus no winning proof tree for \([r_1]\) exists under
that restriction, despite its semantic membership.
Since \(C\) has a one-node proof but its subargument
\([r_1]\) has none, proof-theoretic subargument closure
also fails for that formulation.

In the corrected game, respond with \(C\) to every opponent
defeater of initial \([r_1]\).
Every such defeater contains \(r_3\) and is strictly defeated
by \(C\); no one defeats \(C\).
There are finitely many such arguments.
Thus all opponent branches close with \(C\), and its content
differs from the initial singleton. This is a complete proof,
not merely a repair of one selected branch.
\end{proof}
\end{annexitem}

\begin{annexitem}{annex:p18.is1996-finite}{thm:p18-is1996-finite}{Finite-ground decision procedure}
\begin{proof}
Let there be \(n\) ground rules.
Enumerate their \(2^n\) subsets.
Reject any subset with two different producers of a head.
For each remaining subset, repeatedly select a rule whose
strong premises have already been produced.
If all rules can be selected, their selection order is
a valid argument list.
If the process stops early, a valid list cannot exist:
the first unselected rule in any supposed valid list
would have all its earlier prerequisites selected, contradicting
the stopping condition.
This enumerates precisely the argument contents.

There are finitely many strict extension subsets to consider
for every argument. Validate them by the same ordering
procedure and compute the finite frontiers along their rules.
Enumerate all literal conflict witnesses, compare their
frontiers against the given finite priority relation,
and apply the undercut and empty-incoherence clauses.
Thus attack and defeat are decidable on the finite universe.
Iterate \(F_\Gamma\) until its value stops changing;
at most the number of argument contents can be newly added.
Compute the source's remaining labels from this fixed point
and its attack relation.
The finite theory is finitary, so the corrected proof theorem
decides the same justified arguments and constructs proofs
by descending finite entry ranks.

The two schemata in the body have distinct ground heads
\(p(c),p(f(c)),p(f^2(c)),\ldots\).
Their arguments therefore have unbounded lengths.
This directly shows why finite written schemata cannot
replace the finite-ground hypothesis.
\end{proof}
\end{annexitem}

\begin{annexitem}{annex:p18.is1996-lis-bridge}{thm:p18-is1996-lis-bridge}{The exact fixed-defeat justification bridge}
\begin{proof}
The relation supplied to the fixed-role LIS game is the
relation computed from the present rule theory, without
adding or removing any defeat.
The proof of the semantic theorem established that the
least fixed point using strict defenders equals the
least fixed point using ordinary defenders for this relation.
The latter is precisely the grounded extension used in
Theorem~\ref{thm:p18-lis-grounded-game}.
That theorem identifies it with winning in the corrected
fixed-role game. For finitary \(\Gamma\), the corrected
finite-proof theorem identifies the same set with finite
proof trees in the present construction.
No target, conclusion, or label conversion is needed for
this graph-level assertion.
The changing-role interface asks for additional subtree
and target data, which this argument does not manufacture.
\end{proof}
\end{annexitem}
